\documentclass[11pt]{article}

\usepackage{amsmath,amssymb,amsthm}
\usepackage[margin=1in]{geometry}
\usepackage[colorlinks=true,linkcolor=blue,citecolor=blue,urlcolor=blue]{hyperref}

\title{A Konyagin-Type Lower Bound in an Erd\H{o}s Divisibility Problem}
\author{}
\date{}

\newtheorem{theorem}{Theorem}
\newtheorem{lemma}{Lemma}
\newtheorem{remark}{Remark}

\newcommand{\dd}{\,\mathrm d}
\newcommand{\dist}{\operatorname{dist}}
\newcommand{\ZZ}{\mathbb Z}
\newcommand{\NN}{\mathbb N}

\begin{document}

\maketitle

\section{Introduction}

For a positive integer \(k\), let \(n_k\) be the least integer \(n>2k\) such
that
\[
        p\nmid \prod_{1\le i\le k}(n-i)
\]
for every prime \(p\) with \(k<p<2k\). This is Erd\H{o}s Problem~\#451 in the
Erd\H{o}s Problems website \cite{ErdosProblems451}. Equivalently, \(n_k\)
is the least integer \(n>2k\) for which no prime in \((k,2k)\) divides any of
the \(k\) consecutive integers
\[
        n-k,\ n-k+1,\ldots,\ n-1 .
\]
Such integers exist. Indeed, if
\[
        Q_k=\prod_{k<p<2k}p,
\]
then every sufficiently large multiple \(n\) of \(Q_k\) satisfies
\[
        n-i\equiv -i\not\equiv 0\pmod p
\]
for all \(1\le i\le k\) and all primes \(p\in(k,2k)\).

We shall call an integer \(n>2k\) \emph{\(k\)-admissible} if it satisfies the
above divisibility condition. Thus \(n_k\) is the least \(k\)-admissible
integer.

The purpose of this note is to prove the lower bound
\[
        n_k>
        \exp\!\left(
        c\frac{(\log k)^2}{\log\log k}
        \right)
\]
with an absolute constant \(c>0\).

The idea is simple. We use only primes very close to \(k\). Let
\[
        H=k^{1-\beta},
\]
where \(0<\beta<1\) is fixed. If \(p\in(k,k+H]\) and \(n\) is
\(k\)-admissible, then the least nonnegative residue of \(n\) modulo \(p\) is
either \(0\), or larger than \(k\). Since \(p-k\le H\), this implies
\[
        \left\|\frac np\right\|\le \frac{H}{k}=k^{-\beta},
\]
where \(\|x\|\) denotes the distance from \(x\) to the nearest integer. Thus
every prime \(p\in(k,k+H]\) gives an integer point \(u=p\) near the curve
\[
        v=\frac{n}{u}.
\]
A theorem of Konyagin on integer points close to a smooth curve gives a strong
upper bound for the number of such \(u\), while the Baker--Harman theorem on
primes in short intervals gives a lower bound of order \(H/\log k\). With a
careful choice of parameters these bounds are incompatible when
\[
        n\le
        \exp\!\left(
        c\frac{(\log k)^2}{\log\log k}
        \right).
\]

\section{Main result}

\begin{theorem}
\label{thm:main}
There exists an absolute constant \(c>0\) such that, for all sufficiently large
\(k\),
\[
        n_k>
        \exp\!\left(
        c\frac{(\log k)^2}{\log\log k}
        \right).
\]
Equivalently, for all sufficiently large \(k\), every integer \(n\) satisfying
\[
        2k<n\le
        \exp\!\left(
        c\frac{(\log k)^2}{\log\log k}
        \right)
\]
has the property that
\[
        \prod_{1\le i\le k}(n-i)
\]
is divisible by at least one prime \(p\in(k,2k)\).
\end{theorem}

We first record a preliminary lower bound. The proof is included to make the
argument self-contained.

\begin{lemma}
\label{lem:prelim}
There exists an absolute constant \(c_1>0\) such that, for all sufficiently
large \(k\), every \(k\)-admissible integer \(n\) satisfies
\[
        n\ge c_1\frac{k^2}{\log k}.
\]
\end{lemma}

\begin{proof}
We use the prime number theorem with the classical de la Vall\'ee Poussin error
term; see, for example, \cite[Chapter 15]{Davenport}. In particular, uniformly for \(k\le u<v\le 2k\),
\[
        \pi(v)-\pi(u)
        =
        \int_u^v \frac{\dd t}{\log t}
        +
        O\!\left(k\exp(-c\sqrt{\log k})\right)
\]
for some absolute constant \(c>0\). Hence every interval
\(I\subset(k,2k)\) of length \(|I|\ge k/(\log k)^2\) contains a prime, once
\(k\) is sufficiently large.

Let \(n\) be \(k\)-admissible and put
\[
        T=\frac nk .
\]
We first prove that, for every fixed \(C>0\), no \(k\)-admissible \(n\) can
satisfy
\[
        T\le C\log k
\]
for all sufficiently large \(k\).

We need a simple geometric observation. For every real \(T>2\), there is an
integer \(a\ge1\) such that
\[
        \left(\frac{T-1}{a},\frac{T}{a}\right)\cap(1,2)
\]
has length at least \(1/(4T)\). Indeed, write
\[
        T=m+\alpha,
        \qquad
        m\ge2,
        \quad
        0\le\alpha<1 .
\]
If \(\alpha\ge1/2\), take \(a=m\). Then
\[
        \left(\frac{T-1}{m},\frac{T}{m}\right)\cap(1,2)
        =
        \left(1,1+\frac{\alpha}{m}\right),
\]
whose length is at least \(1/(4T)\). If \(\alpha<1/2\) and \(m=2\), take
\(a=1\); the intersection is \((1+\alpha,2)\), whose length is larger than
\(1/2\). Finally, if \(\alpha<1/2\) and \(m\ge3\), take \(a=m-1\). Then
\[
        \left(\frac{T-1}{m-1},\frac{T}{m-1}\right)
        =
        \left(
        1+\frac{\alpha}{m-1},
        1+\frac{1+\alpha}{m-1}
        \right)
        \subset (1,2),
\]
and its length is \(1/(m-1)\ge1/(4T)\).

Now suppose \(T\le C\log k\). Choose \(a\) as above. Multiplying by \(k\), the
interval
\[
        \left(\frac{n-k}{a},\frac{n}{a}\right)\cap(k,2k)
\]
has length at least
\[
        \frac{k}{4T}
        \ge
        \frac{k}{4C\log k}.
\]
For sufficiently large \(k\), this length is larger than \(k/(\log k)^2\), so
the interval contains a prime \(p\in(k,2k)\). Thus
\[
        \frac{n-k}{a}<p<\frac{n}{a},
\]
and therefore
\[
        n-k<ap<n .
\]
Let
\[
        i=n-ap .
\]
Then \(i\) is an integer satisfying \(0<i<k\), so \(1\le i\le k-1\), and
\[
        p\mid n-i .
\]
This contradicts \(k\)-admissibility. Therefore, for every fixed \(C>0\),
every \(k\)-admissible \(n\) satisfies
\[
        T>C\log k
\]
for all sufficiently large \(k\).

We now strengthen this to \(T\gg k/\log k\). Put
\[
        L_0=(\log k)^2,
        \qquad
        H_0=\left\lfloor\frac{k}{L_0}\right\rfloor,
\]
and let
\[
        \mathcal P=
        \{p\text{ prime}: k<p\le k+H_0\}.
\]
By the prime number theorem,
\[
        |\mathcal P|\ge c_0\frac{k}{L_0\log k}
\]
for some absolute constant \(c_0>0\) and all sufficiently large \(k\). Since
\(H_0<k\), we have
\[
        \mathcal P\subset(k,2k).
\]

Let \(p\in\mathcal P\). Since \(n\) is \(k\)-admissible, the least
nonnegative residue \(r_p\) of \(n\) modulo \(p\) satisfies either \(r_p=0\), or
\(k+1\le r_p\le p-1\). Hence there is an integer \(j_p\) such that
\[
        0\le j_p\le p-k-1\le H_0-1
\]
and
\[
        p\mid n+j_p .
\]
Define
\[
        a_p=\frac{n+j_p}{p} .
\]
Then
\[
        n\le a_pp<n+H_0 .
\]

Let
\[
        \mathcal A=
        \left\{
        a\in\ZZ:
        \frac{n}{k+H_0}\le a<\frac{n+H_0}{k}
        \right\} .
\]
Every \(a_p\) lies in \(\mathcal A\). Writing \(h=H_0/k\), we have
\(h\le1/L_0\), and the length of the interval defining \(\mathcal A\) is
\[
        \frac{n+H_0}{k}-\frac{n}{k+H_0}
        =
        T+h-\frac{T}{1+h}
        =
        \frac{Th}{1+h}+h
        \le
        \frac{T+1}{L_0}.
\]
Thus
\[
        |\mathcal A|\le \frac{T}{L_0}+2 .
\]

For a fixed \(a\in\mathcal A\), the primes \(p\in\mathcal P\) with \(a_p=a\)
satisfy
\[
        n\le ap<n+H_0 .
\]
Hence they lie in an interval of length \(H_0/a\). Since
\[
        a\ge \frac{n}{k+H_0}
        =
        \frac{T}{1+H_0/k}
        \ge
        \frac{T}{2}
\]
for large \(k\), the number of such primes is at most
\[
        1+\frac{H_0}{a}
        \le
        1+\frac{2H_0}{T}.
\]
Therefore
\[
        |\mathcal P|
        \le
        \left(\frac{T}{L_0}+2\right)
        \left(1+\frac{2H_0}{T}\right).
\]
Using \(H_0\le k/L_0\), this gives
\[
        |\mathcal P|
        \le
        \frac{T}{L_0}+2+\frac{2k}{L_0^2}+\frac{4k}{L_0T}.
\]
Combining this with the lower bound for \(|\mathcal P|\), and multiplying by
\(L_0\), we obtain
\[
        c_0\frac{k}{\log k}
        \le
        T+2L_0+\frac{2k}{L_0}+\frac{4k}{T}.
\]

Now take \(C=16/c_0\) in the first part of the proof. For all sufficiently
large \(k\),
\[
        T>C\log k .
\]
Thus
\[
        \frac{4k}{T}
        \le
        \frac{c_0}{4}\frac{k}{\log k}.
\]
Also,
\[
        2L_0+\frac{2k}{L_0}
        =
        o\!\left(\frac{k}{\log k}\right).
\]
Hence, for sufficiently large \(k\),
\[
        2L_0+\frac{2k}{L_0}
        \le
        \frac{c_0}{4}\frac{k}{\log k}.
\]
It follows that
\[
        c_0\frac{k}{\log k}
        \le
        T+\frac{c_0}{2}\frac{k}{\log k},
\]
and so
\[
        T\ge \frac{c_0}{2}\frac{k}{\log k}.
\]
Since \(n=Tk\), the lemma follows with \(c_1=c_0/2\).
\end{proof}

We shall also use primes in short intervals.

\begin{lemma}[Baker--Harman]
\label{lem:short-primes}
There is an absolute constant \(\alpha>0\) such that the following holds. If
\(0<\beta<\alpha\), then, for all sufficiently large \(x\),
\[
        \#\{p\text{ prime}:x<p\le x+x^{1-\beta}\}
        \gg_\beta
        \frac{x^{1-\beta}}{\log x}.
\]
One may take any fixed \(\beta<0.465\).
\end{lemma}

\begin{proof}
We use the Baker--Harman theorem \cite{BakerHarman} in the following standard
form: there is an admissible constant \(\alpha>0\), with \(\alpha=0.465\)
admissible, such that uniformly for sufficiently large \(x\) and
\[
        x^{1-\alpha}\le y\le x,
\]
one has
\[
        \pi(x+y)-\pi(x)\gg \frac{y}{\log x}.
\]
Taking \(y=x^{1-\beta}\), with any fixed \(0<\beta<\alpha\), gives the stated
estimate.
\end{proof}

The main analytic input is Konyagin's integer-point theorem. We state precisely
the form needed below.

\begin{lemma}[Konyagin's integer-point theorem]
\label{lem:konyagin-original}
Let \(r>1\) be an integer, \(N>0\), \(W>1\), and let
\(f\in C^{r+1}[0,N]\). Let \(D_r>0\), \(D_{r+1}>0\), \(\delta>0\), and
\(\lambda\ge1\). Suppose that, for all \(u\in[0,N]\),
\[
        \frac{D_r}{2}\le f^{(r)}(u)\le D_r,
        \qquad
        |f^{(r+1)}(u)|\le D_{r+1} .
\]
Let \(S\subset[0,N]\cap\ZZ\), and suppose that for every \(u\in S\) there
exist \(v\in\ZZ\) and \(w\in\NN\), with \(w\le W\), such that
\[
        \left|f(u)-\frac vw\right|<\delta .
\]
Then
\[
\begin{aligned}
        |S|
        \le C N\bigg[&
        (D_r\lambda^rW^2)^{1/(2r-1)}
        +
        \left(\frac{\delta W^2}{D_r\lambda^r}\right)^{1/(r-1)}
        +
        \left(\frac{D_{r+1}\lambda W^2}{D_r}\right)^{1/(2r)}
        \bigg]
        +2r\lambda,
\end{aligned}
\]
where \(C>0\) is an absolute constant.
\end{lemma}

\begin{proof}
This is Theorem~2, formula~(4), of Konyagin~\cite{Konyagin}. Konyagin states
throughout that all constants denoted by \(c_i\) are absolute; in particular
the constant in this estimate is independent of
\(r,N,W,D_r,D_{r+1},\delta\), and \(\lambda\).
\end{proof}

We shall use only the following specialization to the curve \(A/u\).

\begin{lemma}[Konyagin's estimate for \(A/u\)]
\label{lem:konyagin-special}
Let \(M,L\) be positive integers with \(M\ge2\) and \(1\le L\le M\).
Let \(A>M\), let \(0<\Delta<1/2\), and let \(r\ge2\) be an integer satisfying
\[
        \frac{(r+1)L}{M}\le \log 2 .
\]
Put
\[
        D=\frac{A r!}{M^{r+1}} .
\]
For
\[
        \mathcal R(A;M,L,\Delta)
        :=
        \left\{
        u\in(M,M+L]\cap\ZZ:
        \left\|\frac Au\right\|<\Delta
        \right\},
\]
one has, for every \(\lambda\ge1\),
\[
\begin{aligned}
        |\mathcal R(A;M,L,\Delta)|
        \le C L\bigg[&
        (D\lambda^r)^{1/(2r-1)}
        +
        \left(\frac{\Delta}{D\lambda^r}\right)^{1/(r-1)}
        +
        \left(\frac{(r+1)\lambda}{M}\right)^{1/(2r)}
        \bigg]
        +2r\lambda,
\end{aligned}
\]
where \(C>0\) is an absolute constant.
\end{lemma}

\begin{proof}
Apply Lemma \ref{lem:konyagin-original} on the interval \([0,L]\) to
\[
        f(t)=(-1)^r\frac{A}{M+t} .
\]
Then
\[
        f^{(r)}(t)=\frac{A r!}{(M+t)^{r+1}}>0 .
\]
Since \((r+1)L/M\le\log2\), we have, for \(0\le t\le L\),
\[
        \frac{D}{2}
        \le
        \frac{A r!}{(M+t)^{r+1}}
        \le
        D .
\]
Moreover,
\[
        |f^{(r+1)}(t)|
        =
        \frac{A(r+1)!}{(M+t)^{r+2}}
        \le
        \frac{A(r+1)!}{M^{r+2}}
        =
        \frac{(r+1)D}{M} .
\]
Now let
\[
        S=
        \{t\in[0,L]\cap\ZZ: M+t\in\mathcal R(A;M,L,\Delta)\}.
\]
Since \(M\) is an integer, the map \(t\mapsto M+t\) is a bijection from \(S\)
onto \(\mathcal R(A;M,L,\Delta)\). Hence
\[
        |S|=|\mathcal R(A;M,L,\Delta)|.
\]
For every \(t\in S\), there exists an integer \(m\in\ZZ\) such that
\[
    \left|\frac{A}{M+t}-m\right|<\Delta.
\]
Hence
\[
    |f(t)-(-1)^r m|<\Delta.
\]
Thus the rational-approximation hypothesis of
Lemma~\ref{lem:konyagin-original} is satisfied with
\[
    v=(-1)^r m,\qquad w=1.
\]
Since Lemma~\ref{lem:konyagin-original} is stated for \(W>1\), we may simply
take \(W=2\); the resulting powers of \(W\) are absolute constants and will be
absorbed into the final constant \(C\).

Applying Lemma~\ref{lem:konyagin-original} with
\[
    N=L,\qquad
    D_r=D,\qquad
    D_{r+1}=\frac{(r+1)D}{M},\qquad
    \delta=\Delta,\qquad
    W=2,
\]
and absorbing the harmless powers of \(W=2\) into the absolute constant gives
the stated estimate.
\end{proof}

\section{Proof of the main theorem}

\begin{proof}[Proof of Theorem \ref{thm:main}]
Fix
\[
        \beta=\frac1{10}.
\]
This choice is smaller than \(0.465\), so Lemma \ref{lem:short-primes} applies.
Put
\[
        X=\log k,
        \qquad
        Y=\log X .
\]
Choose a positive constant \(c\) sufficiently small so that
\[
        \frac{\beta}{15c}>20,
        \qquad
        \frac{1}{20c}>20 .
        \tag{C}
\]
We shall prove that, for all sufficiently large \(k\), no \(k\)-admissible
integer \(n\) can satisfy
\[
        2k<n\le
        \exp\!\left(c\frac{X^2}{Y}\right).
        \tag{2.1}
\]

Assume, for contradiction, that such a \(k\)-admissible \(n\) exists for
arbitrarily large \(k\). By Lemma \ref{lem:prelim},
\[
        n\ge c_1\frac{k^2}{X}.
        \tag{2.2}
\]
Let
\[
        H=\lfloor k^{1-\beta}\rfloor,
        \qquad
        \Delta=2k^{-\beta} .
\]
For sufficiently large \(k\), we have \(H<k\), and hence
\[
        (k,k+H]\subset(k,2k).
\]

Let \(p\in(k,k+H]\) be prime, and let \(r_p\) be the least nonnegative residue
of \(n\) modulo \(p\). Since \(n\) is \(k\)-admissible,
\[
        r_p\notin\{1,2,\ldots,k\}.
\]
Thus either \(r_p=0\), or \(k+1\le r_p\le p-1\). If \(r_p=0\), then
\[
        \left\|\frac np\right\|=0 .
\]
If \(r_p>k\), then
\[
        \left\|\frac np\right\|
        \le
        1-\frac{r_p}{p}
        \le
        1-\frac{k+1}{p}
        <
        \frac{p-k}{p}
        \le
        \frac{H}{k} .
\]
Therefore, for every prime \(p\in(k,k+H]\),
\[
        \left\|\frac np\right\|\le \frac{H}{k}<\Delta .
\]
Since \(H=\lfloor k^{1-\beta}\rfloor\), Lemma \ref{lem:short-primes}, applied
with \(x=k\), gives
\[
        \#\{p\text{ prime}:k<p\le k+H\}
        \gg
        \frac{H}{X}.
\]
Hence
\[
        \mathcal N
        :=
        \left|
        \left\{
        u\in(k,k+H]\cap\ZZ:
        \left\|\frac nu\right\|<\Delta
        \right\}
        \right|
        \gg
        \frac{H}{X}.
        \tag{2.3}
\]

We shall now contradict this lower bound by applying Lemma
\ref{lem:konyagin-special}.

For \(j\ge2\), define
\[
        Q_j=\frac{n j!}{k^{j+1}} .
\]
Let \(r\) be the least integer \(j\ge2\) for which
\[
        Q_j\le k^{-\beta}.
        \tag{2.4}
\]
Such an \(r\) exists for all sufficiently large \(k\). Indeed, if
\[
        R=\left\lceil \frac{4cX}{Y}\right\rceil+10,
\]
then, using \((2.1)\) and \(\log(R!)\le R\log R\), we get
\[
        \log Q_R
        \le
        c\frac{X^2}{Y}+R\log R-(R+1)X
        \le
        -2c\frac{X^2}{Y}
\]
for sufficiently large \(k\). In particular \(Q_R\le k^{-\beta}\), once
\(k\) is large. Thus
\[
        2\le r\le R\le 5c\frac{X}{Y}
        \tag{2.5}
\]
for sufficiently large \(k\).

Since \(H/k\le k^{-\beta}\), it follows from \((2.5)\) that
\[
        \frac{(r+1)H}{k}=o(1).
\]
In particular,
\[
        \frac{(r+1)H}{k}\le\log2
\]
for all sufficiently large \(k\). Hence Lemma \ref{lem:konyagin-special} may be
applied with
\[
        M=k,
        \qquad
        L=H,
        \qquad
        A=n,
        \qquad
        D=Q_r,
        \qquad
        \Delta=2k^{-\beta}.
\]

Set
\[
        a_r=\frac{2r-1}{3r-2}
\]
and choose
\[
        \lambda=
        \left(\frac{\Delta^{a_r}}{Q_r}\right)^{1/r} .
        \tag{2.6}
\]
Since \(a_r<1\), \(\Delta=2k^{-\beta}\), and \(Q_r\le k^{-\beta}\), we have
\(\lambda\ge1\) for all sufficiently large \(k\). Lemma
\ref{lem:konyagin-special} gives
\begin{equation}
\begin{aligned}
        \mathcal N
        \le C H\bigg[&
        (Q_r\lambda^r)^{1/(2r-1)}
        +
        \left(\frac{\Delta}{Q_r\lambda^r}\right)^{1/(r-1)}
        +
        \left(\frac{(r+1)\lambda}{k}\right)^{1/(2r)}
        \bigg]
        +2r\lambda .
\end{aligned}
\tag{2.7}
\end{equation}
By the definition of \(\lambda\), the first two terms in the square brackets are
both equal to
\[
        \Delta^{1/(3r-2)} .
\]
Using \((2.5)\) and \((C)\), we get
\[
        \Delta^{1/(3r-2)}
        \ll
        \exp\!\left(-\frac{\beta X}{3r}\right)
        \le
        \exp\!\left(-\frac{\beta Y}{15c}\right)
        =
        X^{-\beta/(15c)}
        \ll
        X^{-20}.
        \tag{2.8}
\]

It remains to estimate the third term in \((2.7)\) and the final term
\(2r\lambda\).

First suppose \(r=2\). Then \(Q_2=2n/k^3\), and \((2.2)\) gives
\[
        Q_2\gg \frac{1}{kX}.
\]
Since \(a_2=3/4\), it follows from \((2.6)\) that
\[
        \lambda
        \ll
        \left(k^{-3\beta/4}\cdot kX\right)^{1/2}
        =
        X^{1/2}k^{1/2-3\beta/8} .
        \tag{2.9}
\]
Therefore
\[
        \left(\frac{3\lambda}{k}\right)^{1/4}
        \ll
        \left(X^{1/2}k^{-1/2-3\beta/8}\right)^{1/4}
        \ll
        k^{-1/10}
        \ll
        X^{-20},
        \tag{2.10}
\]
and also
\[
        r\lambda
        \ll
        X^{1/2}k^{1/2-3\beta/8}
        =
        o\!\left(\frac{H}{X^2}\right),
        \tag{2.11}
\]
because \(H\asymp k^{1-\beta}\) and
\[
        1-\beta-
        \left(\frac12-\frac{3\beta}{8}\right)
        =
        \frac12-\frac{5\beta}{8}
        >0 .
\]

Now suppose \(r\ge3\). By the minimality of \(r\),
\[
        Q_{r-1}=\frac{n(r-1)!}{k^r}>k^{-\beta}.
\]
Since \(Q_r=(r/k)Q_{r-1}\), we have
\[
        Q_r>r k^{-\beta-1}.
\]
Also
\[
        1-a_r=\frac{r-1}{3r-2}<\frac13 .
\]
Hence \((2.6)\) gives
\[
        \lambda^r
        =
        \frac{\Delta^{a_r}}{Q_r}
        \ll
        \frac{k^{-\beta a_r}}{r k^{-\beta-1}}
        \le
        k^{1+\beta/3},
\]
and therefore
\[
        \lambda\ll k^{(1+\beta/3)/r} .
        \tag{2.12}
\]
For \(r\ge3\) and \(\beta=1/10\), the exponent \((1+\beta/3)/r\) is at most
\(31/90\). Since \(r\le5cX/Y\), we have \(r+1\le k^{1/20}\) for sufficiently
large \(k\). Thus, by \((2.12)\),
\[
        \frac{(r+1)\lambda}{k}
        \ll
        k^{-1+1/20+31/90}
        \le
        k^{-1/2}
\]
for all sufficiently large \(k\). Consequently,
\[
        \left(\frac{(r+1)\lambda}{k}\right)^{1/(2r)}
        \ll
        k^{-1/(4r)}
        \le
        \exp\!\left(-\frac{Y}{20c}\right)
        =
        X^{-1/(20c)}
        \ll
        X^{-20},
        \tag{2.13}
\]
where we used \((2.5)\) and \((C)\).

Finally, still assuming \(r\ge3\), \((2.12)\) gives
\[
        r\lambda
        \ll
        r k^{(1+\beta/3)/r}
        \le
        r k^{31/90}.
\]
Since \(H\asymp k^{9/10}\) and \(r\ll X/Y\), this implies
\[
        r\lambda
        =
        o\!\left(\frac{H}{X^2}\right).
        \tag{2.14}
\]

Combining \((2.7)\)--\((2.14)\), we obtain
\[
        \mathcal N
        \ll
        \frac{H}{X^{20}}
        +
        o\!\left(\frac{H}{X^2}\right)
        \ll
        \frac{H}{X^2}.
        \tag{2.15}
\]
But \((2.3)\) gives
\[
    \mathcal N\gg \frac{H}{X}.
\]
Equivalently, there exists an absolute constant \(c_*>0\) such that
\[
    \mathcal N\ge c_*\frac{H}{X}
\]
for all sufficiently large \(k\). On the other hand, \((2.15)\) yields
\[
    \mathcal N\le C_*\frac{H}{X^2}
\]
for some absolute constant \(C_*>0\). Hence
\[
    c_*\le \frac{C_*}{X},
\]
which is impossible for sufficiently large \(k\), since \(X=\log k\to\infty\).
\end{proof}

\begin{remark}
The proof uses only primes \(p\) in the short interval
\[
        (k,k+k^{1-\beta}] .
\]
For each such prime, admissibility forces \(n/p\) to be very close to an
integer. Konyagin's theorem says that this can happen for at most
\(O(H/(\log k)^2)\) integers \(u\in(k,k+H]\) when
\[
        n\le
        \exp\!\left(c\frac{(\log k)^2}{\log\log k}\right),
\]
whereas the Baker--Harman theorem supplies \(\gg H/\log k\) primes in the same
interval. The remaining logarithmic loss comes from the fact that we only get
\(H/\log k\) prime-produced points, not \(H\) points.
\end{remark}

\begin{thebibliography}{9}

\bibitem{BakerHarman}
R. C. Baker and G. Harman,
``The difference between consecutive primes,''
\emph{Proc. London Math. Soc.} (3) \textbf{72} (1996), 261--280.

\bibitem{ErdosProblems451}
T. F. Bloom,
Erd\H{o}s Problem \#451,
\url{https://www.erdosproblems.com/451},
accessed 26 April 2026.

\bibitem{Davenport}
H. Davenport,
\emph{Multiplicative Number Theory},
3rd ed., Graduate Texts in Mathematics, vol. 74,
Springer, 2000.

\bibitem{Konyagin}
S. V. Konyagin,
``Estimates of the least prime factor of a binomial coefficient,''
\emph{Mathematika} \textbf{45} (1998), 41--55.

\end{thebibliography}

\end{document}
